% Mean bibliography entries per NeurIPS paper, 1987-2024.
% Data: ../../data/pgfplots/references_per_year.csv

\begin{figure}[h]
\centering
\begin{tikzpicture}
\begin{axis}[
    width=0.85\linewidth, height=5cm,
    xlabel={Year}, ylabel={Bibliography entries per paper},
    xtick distance=4,
    x tick label style={rotate=45, anchor=east, font=\tiny,
                        /pgf/number format/1000 sep={}},
    ytick align=inside,
    tick label style={font=\tiny},
    label style={font=\small},
    enlarge x limits=0.02,
  ]
  \addplot[color={rgb,255:red,44;green,160;blue,44}, dashdotted, very thick]
    table[x=year, y=neurips, col sep=comma]
    {../../data/pgfplots/references_per_year.csv};
  % ChatGPT launch marker
  \draw[black, dashed, semithick] (axis cs:2022,0) -- (axis cs:2022,100);
\end{axis}
\end{tikzpicture}
\caption{\textbf{Mean bibliography length per NeurIPS paper,
  1987--2025.}
  Reference list length grew from roughly 11 entries per paper
  in 1987 to approximately 63 entries in 2022, a fivefold to
  sixfold increase over 35 years. Mean bibliography length then
  plateaued across 2023 to 2025 (59.4, 60.9, and 59.4 entries).
  Vertical dashed line marks the
  late 2022 availability of instruction-tuned writing
  assistants.}
\label{fig:app-references}
\end{figure}
